\documentclass[../ICIT-Paper-2024.tex]{subfiles}
\usepackage{hyperref} 
\usepackage{tablefootnote}
\usepackage{natbib}
\begin{document}

\section{RELATED WORK}\label{section2}

In this section, we present the methods, research, and models for assessing the credit scores of digital wallets within DeFi. Based on some surveys, 
~\citeauthor{packin2023decentralized}~\cite{packin2023decentralized} identifies two primary directions in the current DeFi credit evaluation: (1) the off-chain integration model and (2) the crypto-native credit score. These approaches have been developed to integrate user information across both TradFi and DeFi. The integration aims to provide a comprehensive assessment of a user's credit score in both Web2 and Web3 spaces~\cite{voshmgir2020token}, thereby blurring the boundaries between DeFi and TradFi.
%Trong phần này, chúng tôi trình bày những phương pháp, nghiên cứu, và mô hình đánh giá điểm tín dụng của ví điện tử trong Defi. Dựa trên những khảo sát của bản thân, Parkin2023 chỉ ra hai hướng phát triển đánh giá tín dụng chính trên Defi hiện nay. Đó là (1) off-chain integration mode và (2) A crypto native credit score. Hai phương pháp này được phát triển nhằm mục đích kết hợp thông tin, dữ liệu người dùng trên cả TraFi và Defi. Việc kết hợp như vậy giúp chúng ta đánh giá toàn diện điểm tín dụng của một người dùng trên cả hai không gian Web2 và Web3. Từ đó xóa bỏ ranh giới giữa Defi và TraFi



\subsection{Off-chain integration model}
(1) In the off-chain integration model, off-chain data (user information from TradFi) are utilized either as a singular data source or combined with on-chain data (DeFi data). Machine learning models 
leverage these data to determine credit scores. All of them (including both data and credit scores) are continually updated, encrypted, and stored publicly on the blockchain. This is intended to enable users to access their data easily and transparently.
%1)In the off-chain integration model, dữ liệu off-chain (thông tin người dùng trên Trafi) được sử dụng như một nguồn dữ liệu duy nhất hoặc được kết hợp với dữ liệu on-chain (dữ liệu Defi). Các mô hình học máy sẽ sử dụng dữ liệu này để determine credit score. Tất cả (bao gồm dữ liệu lẫn credit score) đều được cập nhật liên tục, mã hóa và lưu trữ public trên blockchain. Điều này nhằm mục đích cho phép người dùng có thể tiếp cận dữ liệu của mình một cách dễ dàng và minh bạch.


Identity verification plays a crucial role in constructing an off-chain integration model. Users must provide off-chain data such as personal information, residence, assets, land ownership, etc. This presents a significant challenge because few are willing to disclose such private information to third parties. Moreover, the provision of such data runs counter to the consistent anonymity inherent in DeFi. In our investigations, we have identified the following research articles.
%Để xây dựng một off-chain integration model, Identity verification đóng vai trò quan trọng. Người dùng phải cung cấp off-chain data như các thông tin cá nhân, nơi ở, tài sản, và đất đai,etc. Điều này rất khó khăn vì không ai sẵn sàng cung cấp thông tin private như vậy cho bên thứ ba. Việc cung cấp như vậy cũng đi ngược lại tính ẩn danh nhất quán của Defi. Chúng tôi đã tiến hành khảo sát và tìm được những bài báo nghiên cứu sau.

~\citeauthor{zhu2020blockchain}~\cite{zhu2020blockchain} discusses a blockchain-based approach to identity authentication and intelligent credit reporting. It suggests implementing technologies like multidimensional authentication, multifactor weighted score calculation, and encryption algorithms to secure dynamic identity verification and risk control. It introduces systems for (1) analyzing user behavior, identifying risks in real time, and implementing intelligent risk prevention strategies; (2) leveraging blockchain distributed nodes for collecting credit data and utilizing smart contracts for credit reporting and asset pricing.

~\citeauthor{jain2019blockchain}~\cite{jain2019blockchain} introduces a blockchain-based framework designed to tackle the challenges associated with credit scoring and digital verification, particularly for underprivileged individuals. By establishing a secure digital identity and chronicling past transactions on a blockchain, the framework simplifies the verification process and enhances trust through transparency. The paper's objective is to ensure fair credit access to those lacking fundamental financial amenities and to expedite an efficient digital verification process.

~\citeauthor{patel2020kirti}~\cite{patel2020kirti} proposes KiRTi, a public blockchain-operating deep-learning-based credit-recommender system. KiRTi facilitates intelligent lending between potential borrowers (PB) and potential lenders (PL), eliminating intermediaries like third-party credit-rating agencies. KiRTi employs a long-short-term memory (LSTM) model that generates credit scores (CS) derived from historical transaction data recorded on the blockchain, with smart contracts automating the repayment of loans. Tested on a German credit dataset, KiRTi demonstrated a 97.5\% accuracy rate, with computation and communication costs significantly lower than other cutting-edge methods.

~\citeauthor{cho2021verifiable}~\cite{cho2021verifiable} proposes a model for generating and verifying the revocation of verifiable credentials (VC) in credit-scoring contexts. It incorporates zk-SNARK proofs with cross-revocation verification to establish a trust platform among holders, issuers, and verifiers. The model preserves the decentralized nature of Self-Sovereign Identity (SSI) while allowing for controlled deletion rights, thereby enhancing the processing of credit evaluation data through VCs. The system overcomes certain limitations such as Opt-In, limited revocation rights, and legality issues not addressed by the three preceding systems.

Within the scope of this paper, we will not delve into methods and research addressing these challenges. This is because our approach focuses on a crypto-native credit score. Section~\ref{section2} will elaborate further on this product and research in this field.

\subsection{Crypto native credit score}
\subsubsection{Application Products}
A crypto-native credit score establishes an on-chain identity by leveraging data from various on-chain activities, including loan repayments, trading, and participation in governance. Similar to a conventional credit score tied to an individual's identity and applicable across various platforms, the crypto-native score is linked to the borrower's wallet. As the wallet engages in blockchain activities, its data is gathered, and the credit score is dynamically adjusted ~\cite{packin2023crypto}. ~\citeauthor{packin2023crypto}~\cite{packin2023crypto} introduces pertinent products such as Spectral\footnote{\href{https://docs.spectral.finance/}{https://docs.spectral.finance}}, LedgerScore\footnote{\href{https://www.ledgerscore.com/}{https://www.ledgerscore.com}}, and Prestare\footnote{\href{https://linktr.ee/prestare.finance}{https://linktr.ee/prestare.finance}}, among others.

Spectral introduces an infrastructure tailored for credit risk assessment, offering users the capability to generate a comprehensive Multi-Asset Credit Risk Oracle (MACRO) Score. This score incorporates seven primary factors: DeFi transaction history, liquidation record, loan safety margin, age or time-based metrics, general wallet history, market conditions, and credit mix. These factors are derived from both DeFi and non-DeFi data sources and are processed through sophisticated machine learning algorithms. Similarly, LedgerScore furnishes autonomous crypto credit reports to lenders, verifying borrower identities before evaluating creditworthiness through an analysis of cryptocurrency transactions, income streams, and asset portfolios.

While Spectral and LedgerScore exclusively offer on-chain credit scoring services without directly operating lending protocols, some protocols integrate both functionalities. For instance, Prestare generates an on-chain credit score derived from default risk, which considers the characteristics of the user's assets, alongside historical behavioral risk, assessing the past behavior associated with the user's address. ArcX\footnote{\href{https://arcx.io/}{https://arcx.io}} evaluates creditworthiness by analyzing users' on-chain borrowing activities, incentivizing responsible borrowing behavior. 

RociFi\footnote{\href{https://roci.fi/}{https://roci.fi}}, an under-collateralized lending protocol, mandates borrowers to mint a Non-Fungible Credit Score token, utilizing machine learning algorithms to evaluate credit risk. In cases of default, RociFi publicizes borrower information on social media platforms. Telefy\footnote{\href{https://telefy.finance/\#/}{https://telefy.finance}} serves as both a scoring mechanism and lending platform, gauging creditworthiness based on wallet holdings and transaction history, and offering incentives such as higher scores for holding or staking Tele coins.

Moreover, numerous other Web3 firms actively compete to develop and provide reliable DeFi credit scoring solutions. These include CreDA\footnote{\href{https://www.creda.app/home}{https://www.creda.app}}, Quadrata\footnote{\href{https://quadrata.com/}{https://quadrata.com}}, Credefi\footnote{\href{https://www.credefi.finance/}{https://www.credefi.finance}}, TrueFi\footnote{\href{https://truefi.io/}{https://truefi.io}}, Livesight\footnote{\href{https://www.livesight.com/}{https://www.livesight.com/}}, Bloom\footnote{\href{https://bloom.co/}{https://bloom.co}}, and Masa\footnote{\href{https://www.masa.ai/}{https://www.masa.ai/}}, among others.

All products formulate their score by using features of accounts like “transaction history, liquidation history, amounts owed and repaid, credit mix, and length of credit history”. Some of them also integrate with off-chain data of users. These scores are then scaled to a score range. For example, the score range of Spectral is 300 to 850, and TrueFi is 0 to 255.  At the time of writing, to the best of our knowledge, the remaining competitors do not document their approaches publicly and how they optimize their parameters.

\subsubsection{Research Papers}
Although there have been numerous credit-scoring products for digital wallets, the number of research articles on this matter remains quite modest. As of May 2024, we have found a few articles proposing formulas and methods for calculating credit scores. Here is a survey of our findings.
%Mặc dù đã có nhiều sản phẩm tính điểm tín dụng cho ví điện tử nhưng số lượng bài báo nghiên cứu về vấn đề này còn khá khiêm tốn. Tính đến thời điểm tháng 5 năm 2024, chúng tôi đã tìm được một vài bài báo đề xuất công thức và phương pháp tính điểm tín dụng. Dưới đây là những khảo sát mà chúng tôi đạt được.


~\citeauthor{wolf2022scoring}~\cite{wolf2022scoring} proposes a scoring method for Aave accounts based on parameters such as the account’s age, aggregations of the time series of its historical health factors, interactions with the Aave protocol, the types of assets it borrows and uses as collateral, and more. They endeavor to predict the likelihood that an account—with its current and historical on-chain financial behavior—would fail to repay a subsequent loan. In this context, for Aave accounts, they propose to model the probability that a new position will become eligible for liquidation in its first 90 days. The output will be scaled from [0,1] to [300, 1000] and classified similarly to a FICO score.

The model presented in the article has two limitations. (1) The model is only refined for Aave accounts on Ethereum. For other lending protocols like Compound or Venus, etc on various networks, this model might not be suitable. (2) The model focuses on one parameter, the health factor~\footnote{\href{https://docs.aave.com/faq/borrowing\#what-is-the-health-factor}{https://docs.aave.com}}, which can only be calculated for wallets with outstanding loans, not for wallets that deposit only. This leads to incorrect assessments for wallets that only deposit money.
%Model của bài báo có hai hạn chế. (1) Model chỉ được tinh chỉnh cho aave account trên ethereum. Đối với các lending protocol khác như compound hay venus trên nhiều mạng khác thì model này có thể không phù hợp. (2) Model tập trung vào một thông số là health factor. Thông số này chỉ có thể được tính với những ví có vay còn ví không vay thì không có. Điều này dẫn đến model đánh giá chưa đúng với những ví chỉ gửi tiền.

~\citeauthor{uriawan2024decentralized}~\cite{uriawan2024decentralized} proposes a credit score formula to evaluate the position of a borrower within their lending protocol~\cite{uriawan2022trustlend}. This score will be used by lenders/investors to review and optimize loans. Trustlend proposes the scoring formula as a linear combination of four points. (1) Loan Risk score assesses whether the loan candidate has other loans such as housing, cars, etc. (2) Activity score describes the candidate's occupational activities such as job or business activities. (3) Profile score pertains to the personal data of candidates, including age, education level, and marital status. (4) Social recommendation score is a crucial variable where a borrower receives support directly from other users to add recommendation value.
%Bài báo Trustlend  đề xuất một công thức tính điểm để đánh giá vị thể của một borrower trong lending protocol của họ. ĐIểm này sẽ được các lenders/investors sử dụng để xem xét và tối ưu hóa các khoản vay cho borrowers. Trustlend propose công thức tính điểm là một hàm tuyến tính tổng của 4 loại điểm. (1) Loan Risk score is the component for measure the borrower candidate has the other loan such as housing, car etc. (2) Activity score describing the borrower activity in occupation such as job or business activity. (3) Profile score is the personal data of borrower candidates such as age, education level, and marital status. (4) The social recommendation score is the primary variable the borrower gets support directly from the other users to add recommendation value
~\citeauthor{hartmann2023privacy}~\cite{hartmann2023privacy} proposes calculating a “social score” based on a user's social media data. Their contribution is the development of a fully decentralized lending platform enabled by the Ethereum blockchain that utilizes this social score. This platform aims to provide loan opportunities to users without collateral or sufficient credit scores. The algorithm is based on six hypotheses estimating a user's trustworthiness: (1) Users who link their social media accounts have less to hide. (2) The more a user discloses about themselves, the more trustworthy they appear. (3) Trustworthy users have authentic social media profiles. (4) Larger social networks and higher activity indicate greater creditworthiness. (5) People who make truthful statements are trustworthy.

Both scoring formulas primarily rely on off-chain and do not delve deeply into on-chain data. The omission of on-chain data could make it difficult to accurately assess the financial state of a digital wallet address. Moreover, users are required to authenticate with their digital wallets to contribute off-chain data. This authentication presents a dilemma for Web3 users. Most of them prefer to keep their data anonymous when engaging in DeFi. This issue is similar to the systems within the off-chain integration model. The existing research papers do not offer clear resolutions to this issue.
%Cả hai công thức tính điểm chủ yếu tính dựa trên dữ liệu off-chain web2 và chưa khai thác sâu vào dữ liệu on-chain của web3. Việc bỏ qua dữ liệu web3 sẽ khó có thể đánh giá chính xác tình hình tài chính của một địa chỉ ví điện tử. Bên cạnh đó, người dùng bắt buộc phải xác thực với ví điện tử của mình để cung cấp dữ liệu off-chain. Việc xác thực dẫn đến bất cập cho người dùng web3 khi đa số họ muốn ẩn danh thông tin dữ liệu cá nhân của mình khi tham gia Defi. Bất cập này tương tự như những hệ thống trong off-chain integration model và các bài báo chưa trình bày rõ ràng hướng giải quyết.

Beyond the issues presented, the articles have yet to optimize their scoring parameters according to digital wallet classification. The categorization here is into five groups corresponding to the five rating levels of a FICO Score, which is detailed in Table 1. In this study, we optimize our scoring parameters according to the user classification by the five rating levels. To accomplish this, we need to establish a label set for our digital wallet data. These labels are discussed in detail in section 2.
%Ngoài các vấn đề đã trình bày, các bài báo đều chưa tối ưu tham số điểm của mình theo phân loại ví điện tử. Việc phân loại ví điện tử ở đây là phân loại theo 5 nhóm tương ứng với 5 mức rating của FICO Score. 5 mức rating này được thể hiện chi tiết ở bảng một ở bảng 1. Trong bài báo này, chúng tôi tiến hành tối ưu tham số điểm của mình theo hướng phân loại người dùng theo 5 mức rating. Để làm được điều đó chúng tôi cần xây dựng một tập nhãn cho dữ liệu ví điện tử của mình.

\end{document}